%% SYNTHETIC TEST UNIT generated by writing/tools/make_outline_book.py from FINAL_OUTLINE_v2 -- not book text
\chapter{严格核心 II：深度非弹散射与算符乘积展开}\label{ch:dis}
本单元为测试用的合成文本（由 FINAL\_OUTLINE\_v2 自动生成），陈述内容从略。

\section{第 1 节（测试）}\label{sec:dis-kinematics}

\begin{definition}[DIS 运动学与强子张量]\label{def:dis-kinematics}
（测试用合成陈述：内容从略。）
% uses: none
\end{definition}

\begin{proposition}[强子张量的一般结构 \tagR]\label{prop:dis-W-structure}
（$R_{\mathrm{def}}$：由算符定义直接推出。）
（测试用合成陈述：内容从略。）
非微扰地位：\cref{asm:intro-wightman}。
\uses{def:dis-kinematics, asm:intro-wightman}
\end{proposition}

\begin{definition}[前向 Compton 振幅]\label{def:dis-compton}
（测试用合成陈述：内容从略。）
\uses{def:dis-kinematics}
\end{definition}

\begin{proposition}[解析性与交叉（部分子靶） \tagR]\label{prop:dis-compton-analytic}
（$R_{\mathrm{PT}}$：在重整化微扰论中对红外正规化的部分子态严格。）
（测试用合成陈述：内容从略。）
（本命题只涉及部分子态。）
\uses{def:dis-compton, def:ir-landau-setup, prop:prelim-optical}
\end{proposition}

\section{第 2 节（测试）}\label{sec:dis-compton}

\begin{proposition}[色散关系：Taylor 系数 = 矩（部分子靶） \tagR]\label{prop:dis-dispersion}
（$R_{\mathrm{PT}}$：在重整化微扰论中对红外正规化的部分子态严格。）
（测试用合成陈述：内容从略。）
（本命题只涉及部分子态。）
\uses{prop:dis-compton-analytic, def:dis-compton}
\end{proposition}

\begin{example}[自由夸克的前向振幅]\label{ex:dis-free-quark}
（测试用合成陈述：内容从略。）
\uses{prop:dis-dispersion}
\end{example}

\begin{definition}[扭度与扭度 2 算符]\label{def:dis-twist}
（测试用合成陈述：内容从略。）
\uses{def:prelim-qcd}
\end{definition}

\begin{lemma}[扭度 2 规范不变算符的分类 \tagR]\label{lem:dis-twist-classification}
（$R_{\mathrm{def}}$：由算符定义直接推出。）
（测试用合成陈述：内容从略。）
\uses{def:dis-twist, lem:prelim-boost}
\end{lemma}

\section{第 3 节（测试）}\label{sec:dis-ope}

\begin{lemma}[扭度奇偶性与手征 \tagR]\label{lem:dis-twist-chirality}
（$R_{\mathrm{def}}$：由算符定义直接推出。）
（测试用合成陈述：内容从略。）
\uses{def:dis-twist}
\end{lemma}

\begin{lemma}[树图 OPE \tagR]\label{lem:dis-tree-ope}
（$R_{\mathrm{PT}}$：在重整化微扰论中对红外正规化的部分子态严格。）
（测试用合成陈述：内容从略。）
\uses{def:dis-twist, def:prelim-feynman-rules, def:dis-compton}
\end{lemma}

\begin{proposition}[欧氏非例外动量下 Green 函数的逐图渐近展开 \tagA]\label{prop:dis-ope-euclid}
（测试用合成陈述：内容从略。）
（本命题只涉及部分子态。）
% errtype: E-P
\uses{lem:dis-tree-ope, thm:qcd-renormalizability, thm:ee-euclid-ir, lem:appM-uniqueness, lem:appM-ne}
\end{proposition}
\begin{proofsketch}测试。\end{proofsketch}
\paragraph{地位与缺口}测试：本陈述的地位由\cref{asm:appM-remainder}、\cref{rem:appM-remainder-reduction}决定。

\begin{proposition}[物理矩阵元只需规范不变算符 \tagA]\label{prop:dis-ope-gi}
（测试用合成陈述：内容从略。）
\uses{prop:dis-ope-euclid, thm:prelim-composite-renorm, lem:dis-twist-classification, asm:appM-brst-ope, thm:appM-local-cohomology, lem:appM-alien-vanish, lem:appM-uniqueness, thm:appM-brst-reduction, lem:appM-lbc-formal}
\end{proposition}
\begin{proofsketch}测试。\end{proofsketch}
\paragraph{地位与缺口}测试：本陈述的地位由\cref{asm:appM-brst-ope}、\cref{thm:prelim-composite-renorm}决定。

\section{第 4 节（测试）}\label{sec:dis-np}

\begin{proposition}[前向在壳部分子矩阵元 \tagA]\label{prop:dis-ope-forward}
（测试用合成陈述：内容从略。）
\uses{prop:dis-ope-euclid, prop:dis-ope-gi, lem:prelim-scaleless, asm:appM-forward, lem:appM-uniqueness, cor:appM-forward-ope}
\end{proposition}
\begin{proofsketch}测试。\end{proofsketch}
\paragraph{地位与缺口}测试：本陈述的地位由\cref{asm:appM-forward}、\cref{cor:appM-forward-ope}决定。

\begin{theorem}[重整化微扰论中的欧氏 OPE \tagA]\label{thm:dis-ope}
（测试用合成陈述：内容从略。）
\uses{prop:dis-ope-euclid, prop:dis-ope-gi, prop:dis-ope-forward}
\end{theorem}
\begin{proofsketch}测试。\end{proofsketch}
\paragraph{地位与缺口}测试：本陈述的地位由\cref{asm:appM-brst-ope}、\cref{asm:appM-forward}、\cref{asm:intro-np-ope}、\cref{asm:intro-af-coeff}、\cref{cor:appM-forward-ope}决定。

\begin{proposition}[奇扭度系数的手征压低 \tagR]\label{prop:dis-chiral-suppression}
（$R_{\mathrm{PT}}$：在重整化微扰论中对红外正规化的部分子态严格。）
（测试用合成陈述：内容从略。）
（本命题只涉及部分子态。）
\uses{lem:dis-twist-chirality, def:prelim-feynman-rules}
\end{proposition}

\begin{proposition}[由部分子结构函数提取 Wilson 系数 \tagR]\label{prop:dis-wilson-partons}
（$R_{\mathrm{OPE}}$：以 OPE 为前提而严格。QCD 中含夸克的微扰 OPE 是全阶论证（\cref{thm:dis-ope}），严格证明只对标量理论与纯 Yang--Mills 的规范不变算符给出；对强子态还须非微扰 OPE（\cref{asm:intro-np-ope}）。）
（测试用合成陈述：内容从略。）
（本命题只涉及部分子态。）
\uses{thm:dis-ope, prop:dis-ope-forward, lem:prelim-scaleless, prop:dis-dispersion}
\end{proposition}

\section{第 5 节（测试）}\label{sec:dis-moments}

\begin{proposition}[Wilson 系数与约化矩阵元的重整化群 \tagR]\label{prop:dis-ope-rg}
（$R_{\mathrm{OPE}}$：以 OPE 为前提。）
（测试用合成陈述：内容从略。）
非微扰地位：\cref{asm:intro-wightman}、\cref{asm:intro-composite}。
\uses{prop:dis-ope-euclid, prop:dis-ope-gi, prop:dis-ope-forward, thm:dis-ope, thm:prelim-composite-renorm, prop:appJ-local-operators, prop:qcd-rge, asm:intro-wightman, asm:intro-composite}
\end{proposition}

\begin{lemma}[欧氏点在 $\omega$ 平面上的位置 \tagR]\label{lem:dis-euclid-segment}
（$R_{\mathrm{def}}$：由算符定义直接推出。）
（测试用合成陈述：内容从略。）
\uses{def:dis-compton}
\end{lemma}

\begin{proposition}[由 Wightman 公理得到的有限减除色散关系 \tagR]\label{prop:dis-dispersion-np}
（$R_{\mathrm{cond}}$：在所列编号假设之下证明完整；结论的地位不高于其中最弱的假设。）
前提：\cref{asm:intro-wightman}。
（测试用合成陈述：内容从略。）
\uses{def:dis-compton, prop:dis-W-structure, lem:appK-entire, asm:intro-wightman, lem:appM-forward-tube, lem:appM-hatP-zero, lem:appM-positivity, lem:appM-green}
\end{proposition}

\begin{proposition}[低阶矩与增长条件的等价 \tagR]\label{prop:dis-low-moments}
（$R_{\mathrm{def}}$：由算符定义直接推出。）
（测试用合成陈述：内容从略。）
\uses{def:dis-compton}
\end{proposition}

\section{第 6 节（测试）}\label{sec:dis-rg}

\begin{lemma}[OPE 在线段上的多项式逼近 \tagR]\label{lem:dis-ope-segment}
（$R_{\mathrm{cond}}$：在所列编号假设之下证明完整；结论的地位不高于其中最弱的假设。）
前提：\cref{asm:intro-np-ope}。
（测试用合成陈述：内容从略。）
\uses{lem:dis-euclid-segment, lem:dis-twist-classification, lem:appK-markov, asm:intro-np-ope}
\end{lemma}

\begin{theorem}[由定域 OPE 得到固定 $n$ 的 Taylor 系数 \tagR]\label{thm:dis-fixed-n}
（$R_{\mathrm{cond}}$：在所列编号假设之下证明完整；结论的地位不高于其中最弱的假设。）
前提：\cref{asm:intro-np-ope}、\cref{asm:intro-wightman}、\cref{lem:dis-ope-segment}、\cref{prop:dis-dispersion-np}。
（测试用合成陈述：内容从略。）
\uses{lem:dis-ope-segment, prop:dis-dispersion-np, prop:dis-W-structure, lem:appK-markov, asm:intro-np-ope, asm:intro-wightman}
\end{theorem}

\begin{theorem}[DIS 矩因子化的非微扰条件定理 \tagR]\label{thm:dis-moments-np}
（$R_{\mathrm{cond}}$：在所列编号假设之下证明完整；结论的地位不高于其中最弱的假设。）
前提：\cref{asm:intro-wightman}、\cref{asm:intro-composite}、\cref{asm:intro-np-ope}、\cref{asm:intro-af-coeff}、\cref{asm:intro-growth}。
（测试用合成陈述：内容从略。）
\uses{thm:dis-fixed-n, prop:dis-dispersion-np, prop:dis-low-moments, thm:prelim-composite-renorm, prop:qcd-rge, thm:qcd-asymptotic-freedom, asm:intro-wightman, asm:intro-composite, asm:intro-np-ope, asm:intro-af-coeff, asm:intro-growth, lem:appM-norm-bound}
\end{theorem}

\begin{corollary}[$x$ 空间的弱收敛 \tagR]\label{cor:dis-weak-xspace}
（$R_{\mathrm{cond}}$：在所列编号假设之下证明完整；结论的地位不高于其中最弱的假设。）
前提：\cref{asm:intro-wightman}、\cref{asm:intro-composite}、\cref{asm:intro-np-ope}、\cref{asm:intro-af-coeff}、\cref{thm:dis-moments-np}。
（测试用合成陈述：内容从略。）
\uses{thm:dis-moments-np, lem:appA-weak-convergence, lem:appK-hausdorff, thm:appK-even-moments, asm:intro-wightman, asm:intro-composite, asm:intro-np-ope, asm:intro-af-coeff}
\end{corollary}

\section{第 7 节（测试）}\label{sec:dis-wilson}

\begin{theorem}[DIS 结构函数矩的因子化 \tagR]\label{thm:dis-moments}
（$R_{\mathrm{OPE}}$：以 OPE 为前提。）
（测试用合成陈述：内容从略。）
非微扰地位：\cref{asm:intro-wightman}、\cref{asm:intro-composite}、\cref{asm:intro-np-ope}、\cref{asm:intro-af-coeff}、\cref{asm:intro-growth}。
\uses{def:dis-kinematics, prop:dis-dispersion, prop:dis-ope-forward, thm:dis-ope, thm:dis-fixed-n, prop:dis-dispersion-np, thm:dis-moments-np, asm:intro-wightman, asm:intro-composite, asm:intro-np-ope, asm:intro-af-coeff, asm:intro-growth}
\end{theorem}

\begin{corollary}[标度破坏 \tagR]\label{cor:dis-scaling}
（$R_{\mathrm{OPE}}$：以 OPE 为前提。）
（测试用合成陈述：内容从略。）
非微扰地位：\cref{asm:intro-wightman}、\cref{asm:intro-composite}、\cref{asm:intro-np-ope}、\cref{asm:intro-af-coeff}、\cref{asm:intro-growth}。
\uses{thm:dis-moments, prop:dis-ope-rg, prop:prelim-beta}
\end{corollary}

\begin{example}[$n=2$：能动张量]\label{ex:dis-n2}
（测试用合成陈述：内容从略。）
\uses{thm:dis-moments, prop:dis-ope-rg, prop:dis-low-moments}
\end{example}

\begin{proposition}[单圈 Wilson 系数 \tagR]\label{prop:dis-c-oneloop}
（$R_{\mathrm{PT}}$：在重整化微扰论中对红外正规化的部分子态严格。）
（测试用合成陈述：内容从略。）
（本命题只涉及部分子态。）
\uses{lem:dis-tree-ope, prop:sudakov-oneloop, lem:prelim-plus-expansion, lem:prelim-mellin, prop:dis-dispersion}
\end{proposition}

\section{第 8 节（测试）}\label{sec:dis-xspace}

\begin{example}[单圈 OPE 检验]\label{ex:dis-oneloop-ope}
（测试用合成陈述：内容从略。）
\uses{prop:dis-c-oneloop, prop:dis-ope-forward}
\end{example}

\begin{remark}[偶数矩决定 C 偶组合（数学）]\label{rem:dis-uniqueness}
（测试用合成陈述：内容从略。）
\uses{thm:appK-even-moments, cor:appK-pdf, prop:appK-continuation, thm:appK-carlson}
\end{remark}

\begin{remark}[矩的 $O(1/Q^2)$ 不蕴含逐点 $O(1/Q^2)$]\label{rem:dis-nonuniform}
（测试用合成陈述：内容从略。）
% uses: none
\end{remark}

\begin{proposition}[光锥展开 \tagA]\label{prop:dis-light-cone}
（测试用合成陈述：内容从略。）
\uses{prop:dis-ope-euclid, def:dis-twist}
\end{proposition}
\begin{proofsketch}测试。\end{proofsketch}
\paragraph{地位与缺口}测试：见文献。

\section{第 9 节（测试）}\label{sec:dis-higher-twist}

\begin{proposition}[扭度 4 修正的算符结构 \tagR]\label{prop:dis-twist4}
（$R_{\mathrm{OPE}}$：以 OPE 为前提。）
（测试用合成陈述：内容从略。）
非微扰地位：\cref{asm:intro-wightman}、\cref{asm:intro-composite}、\cref{asm:intro-np-ope}、\cref{asm:intro-af-coeff}。
\uses{thm:dis-ope, thm:prelim-composite-renorm, def:dis-twist, asm:intro-wightman, asm:intro-composite, asm:intro-np-ope, asm:intro-af-coeff}
\end{proposition}

\begin{proposition}[靶质量修正 \tagR]\label{prop:dis-tmc}
（$R_{\mathrm{OPE}}$：以 OPE 为前提。）
（测试用合成陈述：内容从略。）
非微扰地位：\cref{thm:dis-moments}。
\uses{thm:dis-moments, def:dis-twist, asm:intro-wightman, asm:intro-composite, asm:intro-np-ope, asm:intro-af-coeff}
\end{proposition}

\begin{example}[第二阶矩的靶质量修正]\label{ex:dis-tmc-n2}
（测试用合成陈述：内容从略。）
\uses{prop:dis-tmc}
\end{example}

\begin{theorem}[Bjorken 求和规则 \tagR]\label{thm:dis-bjorken}
（$R_{\mathrm{OPE}}$：以 OPE 为前提。）
（测试用合成陈述：内容从略。）
非微扰地位：\cref{asm:intro-wightman}、\cref{asm:intro-composite}、\cref{asm:intro-np-ope}、\cref{asm:intro-af-coeff}、\cref{asm:intro-growth}。
\uses{thm:dis-ope, def:prelim-gamma5, prop:appI-larin, asm:intro-wightman, asm:intro-composite, asm:intro-np-ope, asm:intro-af-coeff, asm:intro-growth}
\end{theorem}

\section{第 10 节（测试）}\label{sec:dis-polarized}

\begin{remark}[Crewther 关系]\label{rem:dis-crewther}
（测试用合成陈述：内容从略。）
% uses: none
\end{remark}
